\subsection{Angular sectors.}

We shall first assume, without any further hypothesis, that E is an oriented plane over the ordered field A. We shall say that three half-lines \(D_{0}\) , \(D_{1}\) , \(D_{2}\) (with origin 0) of E form a direct sequence if, for \(x_{i} \in D_{i}\) , \(x_{i} \neq 0\) (i = 0, 1, 2), at least two of the bivectors \(x_{0} \wedge x_{1}\) , \(x_{1} \wedge x_{2}\) , \(x_{2} \wedge x_{0}\) are \textgreater0; in this case the sequences \(D_{1}\) , \(D_{2}\) , \(D_{0}\) and \(D_{2}\) , \(D_{0}\) , \(D_{1}\) are also direct. It is clear that three half-lines forming a direct sequence are distinct. Given two half-lines \(D_{1}\) , \(D_{2}\) of E, an open (resp. closed) angular sector with origin \(D_{1}\) and endpoint \(D_{2}\) is the set (or, by abuse of language, the union) of half-lines D such that the sequence \(D_{1}\) , D, \(D_{2}\) is direct (resp. such that \(D = D_{1}\) or \(D = D_{2}\) or that the sequence \(D_{1}\) , D, \(D_{2}\) is direct).

PROPOSITION 11. --- Let E be an oriented plane over an ordered field A, D₀ a ray of E, and G the set of rays of E distinct from D₀. The relation

« \(D_{1} = D_{2}\) , or the sequence \(D_{0}\) , \(D_{1}\) , \(D_{2}\) is direct »

between elements \(D_{1}\) , \(D_{2}\) of G is a total order relation in G.

Indeed the axioms of total order relations are trivially verified, except for transitivity. Let \(D_{1}\) , \(D_{2}\) , \(D_{3}\) be three half-lines such that the sequences, \(D_{0}\) , \(D_{1}\) , \(D_{2}\) and \(D_{0}\) , \(D_{2}\) , \(D_{3}\) are direct; we shall show that the sequence \(D_{0}\) , \(D_{1}\) , \(D_{3}\) is direct. To this end, take a vector \(x_{i} \neq 0\) in \(D_{i} (i = 0, 1, 2, 3)\) , choose a bivector e \textgreater{} 0 and set \(x_{i} \wedge x_{j} = a_{ij} e (a_{ij} \in A)\) . Writing \(e = x_{0} \wedge y (y \in E)\) , and taking \((x_{0}, y)\) as a basis of E, one readily verifies the relation

\[
a _ {0 1} a _ {2 3} + a _ {0 2} a _ {3 1} + a _ {0 3} a _ {1 2} = 0.\tag{22}
\]

This being so, if \(a_{01} \leqslant 0\) , one has \(a_{12} > 0\) and \(a_{20} > 0\) (since the first sequence is direct), then \(a_{23} > 0\) and \(a_{30} > 0\) (since \(a_{02} < 0\) and since the second sequence is direct), whence \(a_{13} > 0\) (by virtue of (22)); therefore the sequence \((\mathrm{D}_0, \mathrm{D}_1, \mathrm{D}_3)\) is direct in this case. Suppose now that \(a_{01} > 0\) . If \(a_{30} \leqslant 0\) , one has \(a_{02} > 0\) and \(a_{23} > 0\) (since the second sequence is direct), then \(a_{12} > 0\) (since \(a_{20} < 0\) and since the first sequence is direct), hence \(a_{13}>0\) (by (22)), and the sequence \((\mathrm{D}_{0}, \mathrm{D}_{1}, \mathrm{D}_{3})\) is direct. Finally, it is evidently the same if \(a_{01}>0\) and \(a_{30}>0\) . QED.

COROLLARY. --- Let \(D_{1}\) and \(D_{2}\) be two distinct rays of E. For every ray \(D_{0}\) of E such that the sequence \(D_{0}\) , \(D_{1}\) , \(D_{2}\) is direct, the set of half-lines D of E such that \(D_{1} < D < D_{2}\) (for the total order relation defined by \(D_{0}\) ) is equal to the open angular sector with origin \(D_{1}\) and endpoint \(D_{2}\) .

Indeed, given a half-line \(D_{3}\) we must show that the relations “the sequence \(D_{1}\) , \(D_{3}\) , \(D_{2}\) is direct” and “the sequences \(D_{0}\) , \(D_{1}\) , \(D_{3}\) and \(D_{0}\) , \(D_{3}\) , \(D_{2}\) are direct” are equivalent. To abbreviate, let us denote by (ijk) the relation “the sequence ( \(D_{i}\) , \(D_{j}\) , \(D_{k}\) ) is direct”. By prop. 11, the conjunction of (132) and (120) implies (130); likewise the conjunction of (201) and (213) implies (203); whence one half of our assertion. Conversely, suppose (012), (013), and (032); since the conjunction of (312) and (320) implies (310) (prop. 11), and since (310) and (013) are incompatible, (312) is false; this proves (132) and completes the proof.

By virtue of the preceding corollary, the open (resp. closed) angular sector with origin \(D_{1}\) and endpoint \(D_{2}\) is denoted \(\left[\begin{matrix}\mathrm{D}_{1},\mathrm{D}_{2}\end{matrix}\right]\) (resp. \(\left[\begin{matrix}\mathrm{D}_{1},\mathrm{D}_{2}\end{matrix}\right]\) ), it being understood that these are intervals for the order structure defined by any half-line \(D_{0}\) such that the sequence \(D_{0},D_{1},D_{2}\) is direct.

PROPOSITION 12. — Let A be a maximal ordered field, E an oriented plane over A, D₀ a ray of E, and G the set of rays of E distinct from D₀. The totally ordered sets A and G (prop. 11) are isomorphic.

Let, in fact, \((x, y)\) be a basis of E such that \(x \in -\mathrm{D}_{0}\) and the bivector \(x \wedge y\) be \(> 0\) . To each element \(t\) of A we associate the half-line \(f(t)\) to which the vector \((1 - t^{2})x + 2ty\) belongs. It is clear that \(f(\mathrm{A}) \subset \mathrm{G}\) . We show that \(f\) is strictly increasing. Indeed, for the sequence \(\mathrm{D}_{0}\) , \(f(t)\) , \(f(t')\) ( \(t, t'\) in A) to be direct, it is necessary and sufficient, by definition, that at least two of the elements

\[
- 2 t, \quad (1 - t ^ {2}) 2 t ^ {\prime} - (1 - t ^ {\prime 2}) 2 t, \quad 2 t ^ {\prime}
\]

be \textgreater0. Now the second is equal to \(2(t' - t)(1 + tt')\) . Therefore, if \(t < t'\) , one has either \(tt' \geqslant 0\) , hence \(t' > 0\) or -t \textgreater{} 0, or \(tt' < 0\) , hence -t \textgreater{} 0 and \(t' > 0\) ; in any case \(D_{0}\) , \(f(t)\) , \(f(t')\) is direct. Since A is totally ordered, f is an isomorphism of A onto the ordered set \(f(A)\) (Set Theory, Chap. III, § 1, no. 14, Prop. 13).

It remains to show that \(f\) is surjective. For this, consider the positive form \(\Phi\) on \(E\) such that \((x, y)\) be an orthonormal basis for \(\Phi\) . For every half-line \(D \in G\) , there exists one and only one angle \(\varphi\) such that \(2\varphi = (-D_0, D)\) ( \(n^o 1\) , Prop. 3); since \((-D_0, D)\) is not the straight angle, \(\varphi\) is not the right angle, and \(tg\varphi\) is therefore finite. Then, by virtue of formulas (16) ( \(n^o 3\) ), we have \(D = f(tg\varphi)\) . This completes the proof.

Exercises. --- 1) With the notations of no. 1, set Q(x) = Φ(x, x); the vector space E is canonically identified with C⁻(Q) (§ 9, no. 1). For every z ∈ C⁺(Q) and every x ∈ E, one has zx ∈ E; show that x ↦ zx is an element of A(Φ), and that z ↦ sz is an isomorphism from C⁺(Q) onto the algebra A(Φ).

\begin{enumerate}
\def\labelenumi{\arabic{enumi})}
\setcounter{enumi}{1}
\tightlist
\item
  The hypotheses and notation are those of no. 1 and of Exercise 1.
\end{enumerate}

\begin{enumerate}
\def\labelenumi{\alph{enumi})}
\item
  Let C be the set of \(x \in E\) such that \(\Phi(x, x) = 1\) (“unit circle”), and let \(\mathfrak{D}\) be the set of lines D whose intersection with C is nonempty (and consequently consists of two opposite elements of E). One calls a pointed line any pair \(\Delta = (D, z)\) formed by a line \(D \in \mathfrak{D}\) and one of the points \(z \in D \cap C\) . Show that if \(\Delta_1 = (D_1, z_1)\) , \(\Delta_2 = (D_2, z_2)\) are two pointed lines, there exists one and only one rotation \(u\) such that \(u(z_1) = z_2\) (and consequently \(u(D_1) = D_2\) ), which one expresses by writing \(u(\Delta_1) = \Delta_2\) . In the set of pairs \((\Delta_1, \Delta_2)\) of pointed lines, the relation “there exists a rotation \(u\) such that \(u(\Delta_1) = \Delta_1'\) and \(u(\Delta_2) = \Delta_2'\) ” is an equivalence relation. The set \(\mathfrak{A}_1\) of equivalence classes of pointed lines under this relation is called the set of angles of pointed lines, and the equivalence class to which a pair \((\Delta_1, \Delta_2)\) of pointed lines belongs is called the angle of this pair and is denoted \((\widehat{\Delta_1, \Delta_2})\) ; the relation \((\widehat{\Delta_1, \Delta_2}) = (\widehat{\Delta_1'}, \widehat{\Delta_2'})\) is equivalent to \((\widehat{\Delta_1, \Delta_1'}) = (\widehat{\Delta_2, \Delta_2'})\) and the rotation that maps \(\Delta_1\) to \(\Delta_2\) is called the rotation of angle \(\theta = (\Delta_1, \Delta_2)\) and denoted \(h_1(\theta)\) ; \(h_1\) is a bijection from \(\mathfrak{A}_1\) on \(O^+\) and one transports to \(\mathfrak{A}_1\) by means of \(h_1^{-1}\) the commutative group structure of \(O^+\) , writing the group \(\mathfrak{A}_1\) thus defined additively; one also calls a straight angle in \(\mathfrak{A}_1\) the angle of the pair formed by a pointed line \((D, z)\) and of the "opposite" pointed line ( \(D, -z\) ), which corresponds to the rotation \(x \to -x\) .
\item
  In the set \(\mathfrak{D}_{0} \supset \mathfrak{D}\) of non-isotropic lines, the notion of angle between lines is defined as in no. 3, the group \(\mathfrak{A}_{0}\) (using cor. 3 of prop. 1 of no. 1), the right angle in \(\mathfrak{A}_{0}\) and the canonical bijection \(h\) of \(\mathfrak{A}_{0}\) onto \(S^{+}/H\) . With the notations of the cor. of prop. 3, we set \(\bar{d} = h_{1}^{-1} \circ d \circ h\) and \(\bar{i} = h^{-1} \circ i \circ h_{1}\) so that \(\bar{i}\) is a homomorphism of \(\mathfrak{A}_{1}\) into \(\mathfrak{A}_{0}\) and \(\bar{d}\) a homomorphism of \(\mathfrak{A}_{0}\) into \(\mathfrak{A}_{1}\) ; \(\bar{d}\) is bijective, and the kernel of \(\bar{i}\) consists of 0 and the straight angle; moreover, we have \(\bar{d}(\bar{i}(\theta)) = 2\theta\) for \(\theta \in \mathfrak{A}_{1}\) and \(\bar{i}(\bar{d}(\varphi)) = 2\varphi\) for \(\varphi \in \mathfrak{A}_{0}\) . For \(\bar{i}\) is surjective (in other words, for \(\mathfrak{D}\) to be the set of all non-isotropic lines), it is necessary and sufficient that the field A be Pythagorean (chap. VI, § 2, exerc. 8 d)) and that there exist an orthonormal basis for \(\Phi\) ; this is equivalent to saying that \(\Phi(x, x)\) is a square for every \(x \in E\) .
\end{enumerate}

\begin{enumerate}
\def\labelenumi{\arabic{enumi})}
\setcounter{enumi}{2}
\tightlist
\item
  The hypotheses and notations are those of exerc. 2.
\end{enumerate}

\begin{enumerate}
\def\labelenumi{\alph{enumi})}
\tightlist
\item
  Let \(s\) be a symmetry with respect to a non-isotropic line D ( \(\S 6\) , no. 4). For every pointed line \(\Delta_{1} = (D_{1}, z_{1})\) , let
\end{enumerate}

\[
\Delta_ {2} = s (\Delta_ {1}) = (s (\mathrm{D} _ {1}), s (z _ {1})),
\]

and let \(\varphi = (\widehat{\mathrm{D}_1, \mathrm{D}})\) ; show that one has \((\widehat{\Delta_1, \Delta_2}) = \bar{d}(\varphi)\) .

\begin{enumerate}
\def\labelenumi{\alph{enumi})}
\setcounter{enumi}{1}
\item
  Show that every orthogonal transformation with determinant -1 is a symmetry \(s\) with respect to a non-isotropic line (cf. § 6, exerc. 15 e)); for every rotation \(u\) , one has \(sus^{-1} = u^{-1}\) .
\item
  If \(x, y\) are any two points of E such that \(\Phi(x, x) = \Phi(y, y) \neq 0\) , there exists one and only one symmetry with respect to a non-isotropic line, which transforms \(x\) to \(y\) .
\item
  Let \(s, s'\) be the symmetries with respect to two non-isotropic lines D, D', and let \(\varphi = (\widehat{\mathrm{D},\mathrm{D}'})\) ; for \(s's\) to be a rotation of angle \(\theta\) , it is necessary and sufficient that \(\bar{d}(\varphi) = \theta\) .
\item
  Show that the commutator subgroup of the orthogonal group \(\mathbf{O}(\mathrm{Q})\) is the image of \(\mathfrak{A}_{1}\) under the homomorphism \(\theta \to h_{1}(2\theta)\) ( \(\S 6\) , exerc. \(17a\) ); for this image to be equal to \(\mathrm{O}^{+}\) , it is necessary and sufficient that the homomorphism \(\bar{i}\) be surjective (exerc. \(2b\) ).
\end{enumerate}

\begin{enumerate}
\def\labelenumi{\arabic{enumi})}
\setcounter{enumi}{3}
\item
  The hypotheses and notation are those of exercise 2. Let \(a\) , \(b\) be two points of the circle C, \(\Delta_{a}\) , \(\Delta_{b}\) the pointed lines passing through \(a\) and \(b\) (and by 0) respectively, and let \(\theta = (\widehat{\Delta_{a}, \Delta_{b}})\) . For every \(x \in E\) , distinct from \(a\) and \(b\) , let D \(_{xa}\) (resp. D \(_{xb}\) ) be the affine line passing through \(a\) and \(x\) (resp. \(b\) and \(x\) ), and let D \(_{xa}'\) (resp. D \(_{xb}'\) ) be the direction of D \(_{xa}\) (resp. D \(_{xb}\) ); show that for \(x \in C\) ( \(x\) distinct from \(a\) and \(b\) ), it is necessary and sufficient that the angle \(\varphi = (\widehat{\mathrm{D}_{xa}', \mathrm{D}_{xb}'})\) be such that \(\bar{d}(\varphi) = \theta\) (using exercise 3). How does this result change when \(x = a\) or \(x = b\) (cf. § 6, exerc. 25) \(a\) )) ?
\item
  Using the notation of nos. 1 and 2 and of Exercise 2, we assume that \(\Phi\) has index 1, in other words that \(\delta\) is a square \(\gamma^{2}\) in A; we denote by D \(_{1}\) , D \(_{2}\) the isotropic lines of E, containing respectively the vectors \(e_{1}-\frac{1}{\gamma}e_{2}\) and \(e_{1}+\frac{1}{\gamma}e_{2}\) .
\end{enumerate}

\begin{enumerate}
\def\labelenumi{\alph{enumi})}
\item
  Show that the group of rotations \(O^{+}\) is isomorphic to the multiplicative group A* of the field A.
\item
  For every angle \(\theta \in \mathfrak{A}_1\) , set \(e_w(\theta) = c_w(h_1(\theta)) + \gamma s_w(h_1(\theta))\) . Show that \(\theta \to e_w(\theta)\) is an isomorphism of \(\mathfrak{A}_1\) onto the group \(A^*\) .
\item
  Let D and D′ be any two lines, and let \(\varphi = (\widehat{\mathbf{D},\mathbf{D}^{\prime}})\) . Show that the cross-ratio \(\left[ \begin{array}{cc}\mathrm{D}_{1} & \mathrm{D}_{2}\\ \mathrm{D}^{\prime} & \mathrm{D} \end{array} \right]\) (chap. II, 2nd ed., App. III, exerc. 5) is equal to \(e_w(\overline{d} (\varphi))\) ("Laguerre formula"). (Note that D1 and D2 are invariant under any direct similarity, and using exercise 4 c) of chap. II, 2nd ed., App. III, reduce to the case where D = Ae1.)
\end{enumerate}

\begin{enumerate}
\def\labelenumi{\arabic{enumi})}
\setcounter{enumi}{5}
\tightlist
\item
  We assume that A is an ordered field.
\end{enumerate}

\begin{enumerate}
\def\labelenumi{\alph{enumi})}
\item
  Let \(D_{1}\) , \(D_{2}\) be two non-isotropic rays with origin 0. Show that there exists a direct similarity u such that \(u(\mathrm{D}_{1}) = \mathrm{D}_{2}\) ; any other direct similarity having this property is of the form \(\varrho = su\) where s is a homothety with ratio \textgreater{} 0.
\item
  In the set of pairs \((\mathrm{D}_{1}, \mathrm{D}_{2})\) of half-lines that are not isotropic, the relation “there exists a direct similarity u such that \(u(\mathrm{D}_{1}) = \mathrm{D}_{1}'\) and \(u(\mathrm{D}_{2}) = \mathrm{D}_{2}'\) ” is an equivalence relation. The set A of equivalence classes of non-isotropic rays, under this relation, is called the set of angles of (non-isotropic) rays, and the equivalence class of a pair \((\mathrm{D}_{1}, \mathrm{D}_{2})\) of such half-lines is called the angle of this pair and denoted \((\widehat{\mathrm{D}_{1}, \mathrm{D}_{2}})\) ; if \(\theta = (\widehat{\mathrm{D}_{1}, \mathrm{D}_{2}})\) , one says that \(\theta\) is the angle of every direct similarity mapping \(D_{1}\) to \(D_{2}\) ; let \(h_{2}(\theta)\) be the class mod. \(H^{+}\) of these similarities, so that \(h_{2}\) is a bijection from the set A onto the group \(S^{+}/H^{+}\) ; one transports to A by means of \(h_{2}^{-1}\) the commutative group structure of \(S^{+}/H^{+}\) , writing the group A thus defined additively. Define a canonical injection \(\bar{j}\) of the group \(A_{1}\) of angles of pointed lines (exerc. 2) into the group A, such that \(h_{2} \circ \bar{j} \circ h_{2}^{-1}\) is the canonical injection j of \(O^{+}\) into \(S^{+}/H^{+}\) . For \(\bar{j}\) to be surjective, it is necessary and sufficient that the homomorphism \(\bar{i}\) of \(A_{1}\) into \(A_{0}\) (exerc. 2 b)) be surjective.
\item
  Show that in \(\mathfrak{A}\) the equation \(2\theta = 0\) has 2 solutions if \(\delta < 0\) and 4 solutions if \(\delta > 0\) . In the first case, the solution \(\varpi \neq 0\) of this equation is also called the straight angle.
\end{enumerate}

¶ 7) Under the hypotheses and notation of exercise 6, suppose the homomorphism \(\bar{j}\) bijective; we then define cos θ and sin θ for every θ ∈ A as in no. 3. Let T be the set of θ ∈ A such that sin θ ≥ 0.

\begin{enumerate}
\def\labelenumi{\alph{enumi})}
\item
  Show that for every \(\theta \in T\) , there exists one and only one angle \(\theta' \in T\) such that \(2\theta' = \theta\) ; we set \(\theta' = \theta/2\) .
\item
  Let L be the Z-module of formal linear combinations of the elements of T with coefficients in Z (chap. II, § 1, no. 8); we denote by \(\xi + \eta\) and \(\dot{\cdot} \xi\) the sum and the opposite in L. Let N be the submodule of L generated by the elements of L of the form \(\xi + \eta \dot{\cdot} (\xi + \eta)\) for all pairs \((\xi, \eta)\) of elements of T such that \(\xi + \eta \in T\) (sum taken in the group \(\mathfrak{A}\) ). Let \(f\) be the homomorphism from L to \(\mathfrak{A}\) which extends the canonical injection of T into \(\mathfrak{A}\) , and \(g\) be the endomorphism of L which extends the map \(\theta \to \theta/2\) from T into itself. We have \(f(N) = \{0\}\) and \(g(N) \subset N\) ; by passing to quotients, we deduce from \(f\) a homomorphism \(f\) from M = L/N into \(\mathfrak{A}\) , and from \(g\) an endomorphism \(g\) of M; we set \(g(\mu) = \mu/2\) and if \(g^m\)
\end{enumerate}

is the m-th iterate of g, \(g^{m}(\mu) = 2^{-m}\mu\) ; we have \(2^{m}(2^{-m}\mu) = \mu\) for all \(\mu \in M\) .

\begin{enumerate}
\def\labelenumi{\alph{enumi})}
\setcounter{enumi}{2}
\item
  Show that the restriction to T of the canonical map \(\psi\) from L onto M = L/N is injective, which allows one to identify T with a subset of M by means of \(\psi\) . Show that, if \(\lambda_1, \ldots, \lambda_m\) are elements \(\neq 0\) of T, the sum \(\lambda_1 + \lambda_2 + \cdots + \lambda_m\) cannot be 0 in M (consider the element \(2^{-m}(\lambda_1 + \cdots + \lambda_m)\) and reason by induction on m, noting that these elements belong to T).
\item
  Let \(M_{+}\) be the set of finite sums (in M) of elements of T; show that \(\mathbf{M}_{+}\cap(-\mathbf{M}_{+})=\left\{0\right\}\) and \(\mathbf{M}=\mathbf{M}_{+}\cup(-\mathbf{M}_{+})\) and consequently that \(M_{+}\) is the set of elements \(\geqslant0\) for a totally ordered group structure on M (note that for every \(\mu\in M_{+}\) , there exists an integer m such that \(2^{-m}\mu\in T\) ); this totally ordered group is said to be the group of measures of angles of half-lines. Show that the homomorphism f from M to A is surjective, and that its kernel is the set of integer multiples of \(2\varpi\) , where \(\varpi\) is the straight angle (exerc. 6 c)). Prove that T is identified with the interval \([0,\varpi]\) in M (prove by induction on m that if \(\mu,\lambda_{1},\ldots,\lambda_{m}\) belong to T and if we have \(\lambda_{1}+\cdots+\lambda_{m}\leqslant\mu\) , then \(\lambda_{1}+\cdots+\lambda_{m}\in T\) ). Show that in T (thus identified with an interval of M), the function \(\theta\to\cos\theta\) is strictly decreasing.
\item
  For the totally ordered group M to be Archimedean (chap. VI, § 1, exerc. 31), it is necessary and sufficient that the additive group of the field A be Archimedean. (To see that the condition is necessary, note that if sin θ is infinitely small with respect to the subfield Q of A (chap. VI, § 2, exerc. 1), then so is sin nθ for every integer n. To see that the condition is sufficient, note that if 0 ≤ θ ≤ π/4, then sin 2θ ≥ √2 sin θ).
\end{enumerate}

\begin{enumerate}
\def\labelenumi{\arabic{enumi})}
\setcounter{enumi}{7}
\tightlist
\item
  Let E be an oriented plane over an ordered field A. Let D', D'' be two distinct half-lines; let \(x'\) (resp. \(x''\) ) be a vector \(\neq 0\) in D′ (resp. D″); the angular sector (open or closed) with origin D' and endpoint D'' is said to be salient (resp. reentrant, flat) if \(x' \wedge x'' > 0\) (resp. \(x' \wedge x'' < 0\) , \(x' \wedge x'' = 0\) ).
\end{enumerate}

\begin{enumerate}
\def\labelenumi{\alph{enumi})}
\item
  For there to exist an automorphism of the vector space E mapping an open (resp. closed) angular sector \(\Sigma_{1}\) into an open (resp. closed) angular sector \(\Sigma_{2}\) , it is necessary and sufficient that \(\Sigma_{1}\) and \(\Sigma_{2}\) be both salient, or both reentrant, or both flat.
\item
  Show that the ordered set \(\left[D', D''\right]\) is isomorphic to the interval \(\left[0,1\right]\) of A (first consider the case of a salient sector and note that, in A, any two closed bounded intervals are isomorphic ordered sets).
\item
  With the notations and hypotheses of Exercise 7, define a canonical bijective map from the set T onto a flat angular sector, and show that this map is an isomorphism for the order structures.
\end{enumerate}

\begin{enumerate}
\def\labelenumi{\arabic{enumi})}
\setcounter{enumi}{8}
\item
  Let A be a Pythagorean ordered field, E a finite-dimensional vector space over A, Q a positive nondegenerate quadratic form on E, for which E admits an orthonormal basis (in other words, Q(x) is a square in A for every \(x \in E\) ). Show that the commutator subgroup \(\Omega(Q)\) of the orthogonal group \(\mathbf{O}(Q)\) is the rotation group \(\mathbf{O}^{+}(Q) = \mathbf{S}\mathbf{O}(Q)\) (use Exercise 3 e) of § 10 and Exercise 17 a) of § 6).
\item
  Let A be a maximal ordered field, E a finite-dimensional vector space over A, Q a positive nondegenerate quadratic form on E. For every orthogonal transformation \(u \in \mathbf{O}(Q)\) show that there exists a decomposition of E as a direct sum of pairwise orthogonal subspaces P, N, \(R_i\) ( \(1 \leqslant i \leqslant r\) ) having the following properties: \(1^\circ u(x) = x\) in P, \(u(x) = -x\) in N; \(2^\circ\) each of the \(R_i\) is of dimension 2, we have \(u(R_i) = R_i\) and the restriction of \(u\) to \(R_i\) is a rotation through the angle \(\theta_i\) distinct from 0 and from the straight angle. Moreover, for two decompositions of this nature, the subspaces P and N are the same, as is the sequence of elements \(\cos \theta_i\) up to order (cf. § 7, no. 3, cor. 2 of th. 2). Deduce that every rotation \(u \in \mathbf{O}^+(Q)\) is a commutator \(tst^{-1}s^{-1}\) , where \(s\) and \(t\) are in \(\mathbf{O}(Q)\) and \(s^2 = 1\) (cf. exerc. 3 d)).
\item
  Let A be a maximal ordered field, L the field of quaternions over A (relative to the pair \((-1,-1)\) ), E the subspace (of dimension 3) of L formed by the pure quaternions ( \(\S 9\) , exerc. 15 a)); every quaternion can therefore be written in a unique way \(s=\alpha.1+\nu\) , where \(\alpha\in A\) and \(\nu\in E\) ; we have \(\alpha^{2}-\nu^{2}=\rho.1\) , where \(\rho\in A\) and \(\rho\geqslant0\) ; we set \(\|\nu\|=\sqrt{\rho}\) .
\end{enumerate}

\begin{enumerate}
\def\labelenumi{\alph{enumi})}
\item
  Let \(\varphi(s)\) be the rotation \(x \to sxs^{-1}\) in E, for the positive nondegenerate quadratic form \(x \to \|x\|^2\) on E (§ 9, no. 5, th. 4 and exerc. 15). Show that if \(\nu \neq 0\) , the vectors of the line D ⊂ E containing \(\nu\) are invariant under \(\varphi(s)\) ; the restriction of \(\varphi(s)\) to the orthogonal plane P = D⁰ is a rotation of angle \(\theta\) such that (for a suitable orientation of P) one has \(\operatorname{tg} \frac{\theta}{2} = \| \nu \| / \alpha\) . (If (1, i, j, k) is the canonical basis of L over A, reduce to the case where \(\nu = \beta i\) , \(\beta \in A\) , and then compute \(sjs^{-1}\) ).
\item
  Show that every quaternion of norm 1 can be written \(tst^{-1}s^{-1}\) , where \(s \in L\) , \(t \in E\) (cf. exerc. 10), and that, if \(a\) , \(b\) are two pure quaternions of norm 1, there exists a quaternion \(s\) such that \(b = sas^{-1}\) .
\item
  For two quaternions \(s = \alpha + \nu\) , \(t = \beta + w\) , where \(\alpha, \beta\) are in A, \(\nu, w\) in E, to be permutable, it is necessary and sufficient that \(\nu = \lambda w\) , \(\lambda \in A\) ; for \(st = -ts\) , it is necessary and sufficient that \(\alpha = \beta = 0\) and that the vectors \(\nu\) and \(w\) in E be orthogonal (cf. § 3, exerc. 10).
\end{enumerate}

\begin{enumerate}
\def\labelenumi{\arabic{enumi})}
\setcounter{enumi}{11}
\tightlist
\item
  Let A be a maximal ordered field, and let L be an n-dimensional Euclidean space over A, whose metric form \(\Phi\) is positive nondegenerate; we denote by \(d(x, y)\) the distance \(\sqrt{\Phi(x - y, x - y)}\) between two points of L ( \(\S 7\) , no. 1, Remark). For every point \(c \in L\) and every element \(\rho > 0\) of A, one calls sphere with center c and radius \(\rho\) the set of \(x \in L\) such that \(d(x, c) = \rho\) ; the spheres are therefore nondegenerate affine quadrics in L, admitting a center ( \(\S 6\) , Exercise 25).
\end{enumerate}

\begin{enumerate}
\def\labelenumi{\alph{enumi})}
\item
  Show that at every point x of a sphere S with center c, the tangent hyperplane to S at x (§ 6, Exercise 25) is perpendicular (§ 6, Exercise 22) to the line passing through c and x.
\item
  Let S be a sphere with center c and radius \(\rho\) , let a be a point of L, and let D be a line passing through a and meeting S in two distinct points \(x_{1}\) , \(x_{2}\) (resp. tangent to S at a point x). Show that one has
\end{enumerate}

\[
\Phi (x _ {1} - a, x _ {2} - a) = (d (a, c)) ^ {2} - \rho^ {2} \quad (\text { resp. } (d (x, a)) ^ {2} = (d (a, c)) ^ {2} - \rho^ {2}).
\]

The element \((d(a, c))^{2} - \rho^{2}\) of A is called the power of a with respect to S.

\begin{enumerate}
\def\labelenumi{\alph{enumi})}
\setcounter{enumi}{2}
\item
  Let \(S_{1}\) , \(S_{2}\) be two spheres not having the same center; show that the set of points of L whose powers with respect to \(S_{1}\) and \(S_{2}\) are equal, is a hyperplane perpendicular to the line passing through the centers of the two spheres, and containing their intersection \(S_{1} \cap S_{2}\) ; this hyperplane is said to be the radical hyperplane of \(S_{1}\) and \(S_{2}\) .
\item
  Let \(S_{1}\) , \(S_{2}\) be two spheres with respective centers \(c_{1}\) , \(c_{2}\) and respective radii \(\rho_{1}\) , \(\rho_{2}\) . Show that the following properties are equivalent:
\end{enumerate}

α) The intersection \(S_{1} \cap S_{2}\) is nonempty and, for every point of this intersection, the tangent hyperplanes to \(S_{1}\) and \(S_{2}\) at this point are perpendicular.

\(\beta_{1}\) The power of \(c_{1}\) with respect to \(S_{2}\) is \(\rho_{1}^{2}\) .

\(\beta_{2})\) The power of \(c_{2}\) with respect to \(\mathbf{S}_{1}\) is \(\rho_{2}^{2}\) .

\(\gamma_{1})\) The radical hyperplane of \(\mathbf{S}_{1}\) and \(\mathbf{S}_{2}\) is the polar hyperplane (§ 6, exerc. 25) of \(c_{1}\) with respect to \(\mathbf{S}_{2}\) .

\(\gamma_{2})\) The radical hyperplane of \(\mathbf{S}_{1}\) and \(\mathbf{S}_{2}\) is the polar hyperplane of \(c_{2}\) with respect to \(\mathbf{S}_{1}\) .

When these properties hold, the spheres \(S_{1}\) , \(S_{2}\) are said to be orthogonal.

\begin{enumerate}
\def\labelenumi{\alph{enumi})}
\setcounter{enumi}{4}
\tightlist
\item
  Let \(S_{1}\) , \(S_{2}\) be two orthogonal spheres, \(c_{1}\) , \(c_{2}\) their centers. Show that if \(\varpi_{1}\) , \(\varpi_{2}\) are the powers of a point x with respect to \(S_{1}\) , \(S_{2}\) respectively, we have \(\varpi_{1} + \varpi_{2} = 2\Phi(x - c_{1}, x - c_{2})\) . Converse.
\end{enumerate}

\begin{enumerate}
\def\labelenumi{\arabic{enumi})}
\setcounter{enumi}{12}
\tightlist
\item
  The hypotheses and notation are the same as in Exercise 12. Given a point \(c \in \mathrm{L}\) and an element \(\alpha \neq 0\) of \(\mathbf{A}\) , we call inversion with pole \(c\) and of power \(\alpha\) the involutive permutation \(u\) of the set \(\mathrm{L} - \{c\}\) which is such that, for all \(x \in \mathrm{L} - \{c\}\) , \(u(x)\) belongs to the line passing through \(c\) and \(x\) and satisfies the relation \(\Phi(x - c, u(x) - c) = \alpha\) . By abuse of language, we say that \(u\) is an inversion in \(\mathrm{L}\) .
\end{enumerate}

\begin{enumerate}
\def\labelenumi{\alph{enumi})}
\item
  If \(u\) , \(\nu\) are two inversions with the same pole \(c\) and powers \(\alpha\) , \(\beta\) , \(uv^{-1}\) is the restriction to L - \(\{c'\}\) of the homothety (chap. II, 2nd ed., App. II, exerc. 6) with center \(c\) and ratio \(\alpha\beta^{-1}\) .
\item
  Let S be a sphere (exerc. 12) containing c. Show that the image of S - \{c\} under an inversion with pole c is a hyperplane perpendicular to the line joining c to the center of S (by an abuse of language, this hyperplane is said to be the image of S under the inversion in question). Converse.
\item
  Let S be a sphere not containing c. Show that, if \(\varpi\) is the power of c with respect to S (exerc. 12 b)), the image of S under an inversion with pole c and power \(\alpha\) is the image of S under a homothety with center c and ratio \(\alpha/\varpi\) . If n = 2 and if, for every \(x \in S\) , we denote by T (resp. T') the tangent to S (resp. \(u(S)\) ) at the point x (resp. \(u(x)\) ), by D the line passing through c, x and \(u(x)\) , show that one has \((\widehat{\mathrm{D},\mathrm{T}}) = (\widehat{\mathrm{T}^{\prime},\mathrm{D}})\) .
\item
  Let \(S_{1}\) , \(S_{2}\) be two orthogonal spheres (exerc. 12 d)), \(S_{1}'\) , \(S_{2}'\) their images under an inversion with pole c. If c does not belong to \(S_{1}\) nor to \(S_{2}\) , show that \(S_{1}'\) and \(S_{2}'\) are orthogonal spheres. If \(c \in S_{1}\) and \(c \notin S_{2}\) , \(S_{1}'\) is a hyperplane containing the center of \(S_{2}'\) . If \(c \in S_{1} \cap S_{2}\) , \(S_{1}'\) and \(S_{2}'\) are perpendicular hyperplanes (§ 6, exerc. 22). Converse.
\item
  Let \(u\) be an inversion with pole \(c\) and of power \(\alpha = \rho^2 > 0\) and C the sphere with center \(c\) and radius \(\rho\) . If \(x_1, x_2\) are two distinct points lying on a line passing through \(c\) , and distinct from \(c\) , the following properties are equivalent: \(\alpha\) \(x_1\) and \(x_2\) are transformed into each other by \(u\) ; \(\beta\) \(x_1\) and \(x_2\) are conjugate with respect to C (§ 6, exerc. 25); \(\gamma\) every sphere containing \(x_1\) and \(x_2\) is orthogonal to C. One also says that \(u\) is the inversion with sphere C.
\end{enumerate}

¶ 14) The hypotheses and notations being the same as in exerc. 12, one takes an origin 0 in L. Let E \(_{1}\) be the vector space direct sum of L and a space Af \(_{1}\) of dimension 1; we denote by Q \(_{1}\) the quadratic form on E \(_{1}\) such that for x ∈ L and η ∈ A, we have

\[
\mathrm{Q} _ {1} (x + \eta f _ {1}) = \mathrm{Q} (x) + \eta^ {2},
\]

form which is positive and nondegenerate; we denote by C the sphere with center 0 and radius 1 in E \(_{1}\) (for Q \(_{1}\) ).

In the Euclidean space E \(_{1}\) let s be the inversion with pole -f \(_{1}\) and power 2 (exerc. 13); its restriction s \(_{0}\) to L maps L onto C -\{-f \(_{1}\) \}; s \(_{0}\) (resp. s \(_{0}^{-1}\) ) is called, by abuse of language, the stereographic projection of L onto C (resp. of C onto L) with viewpoint -f \(_{1}\) . For every inversion u in L, with pole c, s \(_{0}\) us \(_{0}^{-1}\) is an involutive permutation of the complement in C of the set \{s \(_{0}\) (c), -f \(_{1}\) \}; we extend it to an involutive permutation u' of C by setting u'(s \(_{0}\) (c)) = -f \(_{1}\) , u'(-f \(_{1}\) ) = s \(_{0}\) (c), and one says that u' is an inversion in C. Likewise, for every symmetry v in L with respect to a hyperplane, s \(_{0}\) vs \(_{0}^{-1}\) is an involutive permutation of the complement in C of the set \{-f \(_{1}\) \}; one extends it to an involutive permutation v' of C by setting v'(-f \(_{1}\) ) = -f \(_{1}\) and we say that v' is a symmetry in C. The group of permutations of C generated by inversions and symmetries is called the conformal group of C (or of L, by abuse of language).

\begin{enumerate}
\def\labelenumi{\alph{enumi})}
\item
  Show that the conformal group of C is generated by the symmetries \(v'\) and the inversions \(u'\) corresponding to the inversions \(u\) in L, of power \(>0\) . (Use exerc. 13 a) and note that in L every translation, as well as the homothety \(x \to -x\) , are products of symmetries with respect to hyperplanes).
\item
  Let u be an inversion in L of power \textgreater{} 0; show that the corresponding inversion \(u'\) in C is the restriction to C of a well-determined transformation \(u_{1}'\) which is either an inversion of power \textgreater{} 0 in \(E_{1}\) whose sphere (exerc. 13 c)) is orthogonal to C, or a symmetry with respect to a hyperplane of \(E_{1}\) passing through 0 (consider in \(E_{1}\) the inversion \(u_{1}\) with the same pole and the same power as u). Formulate the corresponding proposition for the symmetry \(v'\) in C corresponding to a symmetry v in L with respect to a hyperplane.
\item
  In the vector space \(E_{2} = A \times E_{1}\) , consider the quadratic form \(Q_{2}\) such that, for \(\zeta \in A\) , \(y \in E_{1}\) one has
\end{enumerate}

\[
\mathrm{Q} _ {2} ((\zeta , y)) = \zeta^ {2} - \mathrm{Q} _ {1} (y),
\]

a form which is nondegenerate and of signature \((1, n + 1)\) . Identify \(E_{1}\) with its canonical image in projective space \(\mathbf{P}(E_{2})\) (chap. II, 2nd ed., App. III, no. 4). Show (with the notations of b)) that \(u'\) is also the restriction to C of a projective linear map \(\overline{u}''\) obtained by passing to quotients from a transformation \(u''\) of the orthogonal group \(\mathbf{O}(Q_{2})\) , which is a symmetry with respect to a non-isotropic hyperplane in \(E_{2}\) . Formulate the corresponding proposition for \(\varphi'\) . Deduce from this that the conformal group of L is isomorphic to the quotient of the group \(\mathbf{O}(Q_{2})\) by its center (use prop. 5 and exerc. 17 c) of § 6). Conclude from this that every transformation of the conformal group is a product of at most \(n + 2\) transformations which are inversions or symmetries in L (cf. § 6, exerc. 15 e)).

\begin{enumerate}
\def\labelenumi{\alph{enumi})}
\setcounter{enumi}{3}
\tightlist
\item
  Let \(\Sigma\) be the set whose elements are the spheres and the hyperplanes in the affine space L. Deduce from \(b\) ) that there exists a bijection of \(\Sigma\) onto the complement, in \(\mathbf{P}(\mathrm{E}_2)\) , of the set of \(x \in \mathrm{E}_1\) such that \(\mathrm{Q}_1(x) \leqslant 1\) , so that to two orthogonal spheres there correspond two points conjugate with respect to C.
\end{enumerate}

\begin{enumerate}
\def\labelenumi{\arabic{enumi})}
\setcounter{enumi}{14}
\item
  Generalize the definitions and results of Exercises 12 to 14 to the case where A is a Pythagorean field and where there exists an orthonormal basis for \(\Phi\) .
\item
  Let A be a commutative field, V a vector space over A of odd dimension \(2r + 1\) , F the product space A × V, \(\Psi\) a nondegenerate alternating form on F; in the projective space P(F), of dimension \(2r + 1\) , we say that the set \(C_{0}\) of lines that are the canonical images of the totally isotropic planes of F (for \(\Psi\) ) is the (projective) linear complex associated with \(\Psi\) .
\end{enumerate}

Assume henceforth that A is maximally ordered; let \(\Phi\) be a nondegenerate positive symmetric form on V. Let D be the line orthogonal to V (for \(\Psi\) ) in F, which is contained in V, and let H be the hyperplane orthogonal to D (for \(\Phi\) ) in V; in F, H is a non-isotropic subspace for \(\Psi\) , whose orthogonal for \(\Psi\) is therefore a plane P supplementary to H and containing D. In the affine space \(E = \{1\} \times V \subset F\) , one also says that the set C of intersections with E of the totally isotropic planes of F (for \(\Psi\) ) not contained in V, is the (affine) linear complex associated with \(\Psi\) ; the line \(\Delta = P \cap E\) is called the axis of the linear complex C (for the Euclidean space structure defined on E by the metric form \(\Phi\) ).

\begin{enumerate}
\def\labelenumi{\alph{enumi})}
\item
  Show that, in E, every translation equal to a direction vector of \(\Delta\) (chap. II, \(2^{\text{e}}\) ed., App. II, no. 3) leaves C invariant (cf. § 4, exerc. 6).
\item
  Let \((e_i)_{0\leqslant i\leqslant 2r}\) be an orthonormal basis of V for \(\Phi\) , such that \(e_0\in \mathrm{D}\) and that one has \(\Psi (e_{2i - 1},e_{2i}) = \rho_{i} > 0\) for \(1\leqslant i\leqslant r,\Psi (e_j,e_k) = 0\) for the pairs of indices that are not of the form \((e_{2i - 1},e_{2i})\) or \((e_{2i},e_{2i - 1})\)
\end{enumerate}

(§ 7, no. 3, prop. 6). One takes in the affine space E an origin \(a \in \Delta\) , and one sets \(\Psi(a, e_0) = \rho_0\) . Let \(x\) be a vector belonging to the plane generated by \(e_{2i-1}\) and \(e_{2i}\) , and let \(y\) be the vector of this same plane such that \(\Phi(x, y) = 0\) , \(\Phi(y, y) = 1\) and \(x \wedge y = \lambda e_{2i-1} \wedge e_{2i}\) with \(\lambda > 0\) . Let \(E_i\) be the affine linear variety of dimension 3 spanned by the points \(a\) , \(a + e_0\) , \(a + e_{2i-1}\) , \(a + e_{2i}\) in E, and let \(R_x\) be the intersection of \(E_i\) and of the affine hyperplane (in E) generated by the lines of C containing the point \(a + x\) . Show that the direction of the lines orthogonal (for \(\Phi\) ) to the plane \(R_x\) in \(E_i\) is a line \(L_x\) of the plane \(Ae_0 + Ay\) , such that if \(\theta = (\widehat{D, L_x})\) one has

\[
\operatorname{tg} \theta = \frac {\rho_ {i}}{\rho_ {0}} \sqrt {\Phi (x , x)}
\]

when the plane \(Ae_{0} + Ay\) is oriented so that \(e_{0} \wedge y\) is positive.

¶ 17) a) Let A be a commutative field, J : \(\xi \to \overline{\xi}\) an involutive automorphism of A, E a vector space of dimension 2 over A, \(\Phi\) a nondegenerate Hermitian sesquilinear form (not alternating) on E, \((e_1, e_2)\) an orthogonal basis of E for \(\Phi\) ( \(\S\) 6, no. 1, th. 1), such that

\[
\Phi (\xi_ {1} e _ {1} + \xi_ {2} e _ {2}, \eta_ {1} e _ {1} + \eta_ {2} e _ {2}) = \alpha \xi_ {1} \bar {\eta} _ {1} + \beta \xi_ {2} \bar {\eta} _ {2}
\]

(α and β belonging to the field K of invariants of J); we set γ = β/α. We identify the point \(\xi_{1}e_{1} + \xi_{2}e_{2} \in E\) with the element \(\xi_{1} + \xi_{2}\rho\) of the ring B defined by the conditions \(\rho^{2} = -\gamma\) , \(\rho\xi = \bar{\xi}\rho\) for \(\xi \in A\) (§ 3, exerc. 4 a)). For every \(x = \xi_{1} + \xi_{2}\rho \in B\) , set \(\tilde{x} = \bar{\xi}_{1} - \xi_{2}\rho\) and N(x) = x \(\tilde{x}\) = \(\tilde{x}\) x, so that \(x \to \tilde{x}\) is an involutive antiautomorphism of the algebra B (over K) and that one has \(\Phi(x, x) = \alpha N(x)\) . Show that every similarity for Φ whose determinant is equal to the multiplier (also called a direct similarity) can be written in one and only one way \(x \to xy\) , where y is a non-isotropic vector of E, and that its multiplier is N(y).

\begin{enumerate}
\def\labelenumi{\alph{enumi})}
\setcounter{enumi}{1}
\item
  Suppose first that J is the identity, hence K = A. If A has characteristic ≠ 2, recover in this way the results of no. 1. Develop the corresponding theory when A has characteristic 2 (cf. § 4, exerc. 14; one will distinguish two cases, according as γ is or is not a square in A).
\item
  We assume \(J \neq 1\) , so that A is a separable quadratic extension of K. Then \(\Phi\) necessarily satisfies condition (T) (§ 4, no. 2 and exerc. 1); if \(\Phi\) has index 0, B is a reflexive field with center K (chap. VIII, § 11, exerc. 4), and if \(\Phi\) has index 1, B is isomorphic to \(\mathbf{M}_2(\mathrm{K})\) .
\end{enumerate}

\(d)\) We assume \(\mathrm{J} \neq 1\) and A has characteristic \(\neq 2\) ; we then have \(\mathrm{A} = \mathrm{K}(\theta)\) with \(\theta^2 = -\delta \in \mathrm{K}\) . If S is the group of direct similarities for \(\Phi\) , H the group of homotheties in E, with ratio \(\neq 0\) and belonging to K (a group isomorphic to \(\mathrm{K}^*\) ), show that the group S/H is isomorphic to the group of rotations \(\mathrm{O}^+(\mathrm{Q})\) , where Q is a nondegenerate quadratic form on a vector space F of dimension 3 over K, such that there exists an orthogonal basis \((f_1, f_2, f_3)\) of F for which one has

\[
\mathrm{Q} \left(\zeta_ {1} f _ {1} + \zeta_ {2} f _ {2} + \zeta_ {3} f _ {3}\right) = \gamma \zeta_ {1} ^ {2} + \delta \zeta_ {2} ^ {2} + \gamma \delta \zeta_ {3} ^ {2}
\]

(cf. § 9, exerc. 15).

¶ 18) a) Let A be a commutative field, \(\xi \to \bar{\xi}\) an involutive automorphism of A, E a vector space of even dimension \(2m\) over A, \(\Phi\) a nondegenerate Hermitian form of index 0 on E, satisfying condition (T), \(\Delta\) the discriminant of \(\Phi\) with respect to a basis of E. Let M(Φ) be the group of multipliers of similarities for \(\Phi\) (§ 4, exerc. 8). Show that M(Φ) is a subgroup of the multiplicative group of elements of A of the form \(\alpha\bar{\alpha} - (-1)^{m}\beta\bar{\beta}\Delta\) . (Argue by induction on \(m\) , using exerc. 17, as well as the following two remarks: \(1^{\circ}\) if \(u\) is a similarity with multiplier \(\mu\) , \(x\) a vector of E, \(y = u(x)\) and \(z = u(y)\) , there exists a unitary transformation \(\nu\) such that \(\nu(y) = y\) and \(\nu(z) = \mu x\) ; \(2^{\circ}\) if \(\alpha\) , \(\beta\) , \(\lambda\) are three elements \(\neq 0\) of A such that there exist \(a\) , \(b\) , \(c\) , \(d\) in A satisfying the conditions \(\lambda = a\bar{a} + \alpha c\bar{c}\) , \(\lambda = b\bar{b} + \beta dd\bar{d}\) , then there exist \(s\) , \(t\) in A satisfying the condition \(\lambda = s\bar{s} - \alpha\beta t\bar{t}\) ).

\begin{enumerate}
\def\labelenumi{\alph{enumi})}
\setcounter{enumi}{1}
\tightlist
\item
  Let K be a maximal ordered field, \(\mathrm{K}_{1} = \mathrm{K}((t_{1}))\) the field of formal power series in one indeterminate \(t_{1}\) , with coefficients in K (chap. IV, § 5, no. 7), \(\mathrm{A} = \mathrm{K}_{1}((t_{2}))\) the field of formal power series in a second indeterminate \(t_{2}\) , with coefficients in \(K_{1}\) . Let E be a vector space of dimension 6 over A, Q a nondegenerate quadratic form on E, such that there exists an orthogonal basis \((e_{i})\) for which one has
\end{enumerate}

\[
\mathrm{Q} (\sum_ {i = 1} ^ {6} \xi_ {i} e _ {i}) = \xi_ {1} ^ {2} + \xi_ {2} ^ {2} + \xi_ {3} ^ {2} + \xi_ {4} ^ {2} + t _ {1} \xi_ {5} ^ {2} + t _ {2} \xi_ {6} ^ {2}.
\]

Show that there exists no similarity for Q with multiplier \(t_{1}t_{2}\) .

\backmatter
